\section{Non-Pilot Recipients and Connected Communication}
\label{sec:pilot-mobile-non-pilot}

Pilot must be useful before every contact has installed Pilot.  The unified mobile
application therefore supports a governed communication path which begins with an
existing private contact, can use an authorised external provider, and can produce a
bounded invitation for a recipient who is not yet part of the Pilot network.  This is
an adoption mechanism, but it is not permitted to weaken identity, consent, delivery
truth or household isolation.

\subsection{Persona-bound email authority}

Google authorisation is bound to the exact Production Private Identity persona.  The
OAuth flow uses PKCE, requests only reviewed Gmail scopes, stores access and refresh
material behind an opaque mother-private vault reference, and exposes neither tokens
nor the vault reference to the phone or television.  At execution time the Gmail
adapter resolves the connected binding for the exact sending persona.  A different
adult on the same mother cannot borrow it.

The execution path is
\[
 \mathsf{draft}\rightarrow\mathsf{exact\ phone\ review}\rightarrow
 \mathsf{persona\ approval}\rightarrow\mathsf{OAuth\ adapter}\rightarrow
 \mathsf{provider\ receipt}.
\]
The adapter creates a deterministic message identifier for reconciliation, calls the
official Gmail send operation, retains only normalized identifiers and hashes, and
does not retain the raw provider response.  Connecting a real Google account and
qualifying token refresh remain deployment operations; the code does not manufacture
an account or credential.

\subsection{Exact scope and state truth}

Every draft records the selected channel, recipient, subject, body and a bounded list
of attachment descriptors.  Each attachment descriptor includes its identifier,
display name, media type, byte size and SHA-256 digest.  The entire scope is included
in one immutable content hash.  A change to any field invalidates the earlier
approval.

Pilot distinguishes the following states:
\begin{enumerate}
  \item \texttt{awaiting\_private\_approval}: prepared locally and not sent;
  \item \texttt{approved\_pending\_adapter}: approved but no provider call made;
  \item \texttt{provider\_accepted}: the provider verified receipt of the send
        request, but human delivery and reading are unknown;
  \item \texttt{delivered}: a separately authenticated provider delivery event was
        recorded, while human reading still remains unknown; and
  \item \texttt{delivery\_unknown\_reconcile\_required}: an attempt may have happened
        but Pilot cannot safely claim or automatically repeat it.
\end{enumerate}
This separation prevents a common but serious false claim: an HTTP success from a
mail or messaging provider is not evidence that a person received, read or accepted
the request.

\subsection{Reply follow-through}

Received-message summaries retain sender, provider freshness and source provenance,
but not the raw body.  From the private phone a person may convert a reply into a
task, a reminder proposal requiring a private time, or a calendar proposal requiring
private dates.  A sender may also schedule a conditional no-reply reminder after a
provider-accepted message.  These transformations create new governed objects; they
do not inherit permission to send, schedule or execute an external effect.

\subsection{Useful invitations}

For a recipient without Pilot, the sender can create a seven-day, high-entropy bearer
invitation for one exact request.  The public projection contains the request title,
bounded summary, intended display name, channel and expiry.  It omits sender persona
identifiers, household data, contacts, tasks, files, memory and provider credentials.
After the recipient creates or activates their own private identity, the trusted phone
can claim the invitation and open that exact pending request.  No unrelated sender or
household records become visible.

The current invitation origin is the customer's own mother endpoint.  Internet-wide
delivery therefore depends on the later customer-controlled remote-access or opaque
relay deployment.  The invitation format and claim boundary are implemented without
pretending that a public relay already exists.

\subsection{Abuse controls and child protection}

Blocking is private to the acting persona and prevents further drafts to that contact.
Reports are recorded locally with an explicit category and never claim that an
external provider has received a report unless a future authorised provider receipt
proves it.  A daily new-recipient ceiling limits unsolicited fan-out.  Child and teen
external communication retains the existing guardian decision followed by the
child's exact sender approval; joining Pilot does not bypass that sequence.

Official SMS and WhatsApp delivery remain integration gates.  Pilot may register only
an authorised provider adapter capable of returning a verifiable provider identifier.
It must not automate consumer WhatsApp Web, scrape sessions, guess credentials or
label an unsupported handoff as delivered.

\subsection{Verification}

Automated tests cover persona-bound OAuth selection, cross-person rejection, exact
attachment hashing, provider-accepted versus delivered state, authenticated delivery
advancement, reply conversion, invitation expiry and claim isolation, blocking,
report truth and guardian controls.  Provider-network field qualification remains
separate because tests do not send real mail, SMS or WhatsApp messages.
